\uuid{H7Wa}
\titre{Étude de suite}
\auteur{bodin ; hanani ; mzari}
\organisation{exo7}
\licence{CC BY-NC-SA 4.0}
\datecreate{2020-01-20}
\niveau{L1}
\module{Analyse}
\chapitre{Suite}
\sousChapitre{Autre}
\difficulte{3}
\theme{Suites, étude de suite}
\exo7qcmid{lille-1.suites.025}

\contenu{
\begin{question}
\qclassification{121.00}
\qtags{motcle=étude de suite}
Soient $(u_n)$ et $(v_n)$ deux suites réelles. On suppose que $(v_n)$ est croissante non majorée et que $\displaystyle v_n < u_n$ pour tout $n\geq 0$. Que peut-on en déduire ?
\begin{answers}
    \bad{$\displaystyle \lim _{n\to +\infty}v_n<\lim _{n\to +\infty}u_n$}
    \good{La suite $(u_n)$ est divergente.}
    \bad{$\displaystyle \lim _{n\to +\infty}v_n\leq u_0$}
    \good{$\displaystyle \lim _{n\to +\infty}u_n=+\infty$}
\end{answers}
\begin{explanations}
La suite $(v_n)$ est croissante non majorée, donc sa limite est $+\infty$. Il en est de même pour la limite de $(u_n)$, c'est-à-dire $(u_n)$ est divergente et sa limite est $+\infty$.
\end{explanations}
\end{question}
}
